\section{Architecture}


\subsection{Single Binary Design}

The explorer binary contains three major subsystems running concurrently:

\begin{enumerate}
  \item \textbf{Indexer goroutines}: One per chain, each writing to its own SQLite database
  \item \textbf{API server}: HTTP server (Base) reading from all chain databases, serving
        \texttt{/v1/explorer/\{chain\}/*} endpoints
  \item \textbf{Replication}: Per-chain WAL streaming to S3 with PQ encryption
\end{enumerate}

The frontend (Next.js static export) is embedded in the binary via Go's \texttt{embed}
directive and served at the root path.

\subsection{Per-Chain Storage Isolation}

Each chain receives an independent directory containing a SQLite database (WAL mode) and
a BadgerDB key-value store:

\begin{lstlisting}[language=bash]
{data_dir}/{chain_slug}/
  query/indexer.db    # SQLite: blocks, txs, tokens, contracts
  kv/                 # BadgerDB: hash->data, height->block
\end{lstlisting}

\textbf{Properties}:
\begin{itemize}
  \item Zero write contention between chains (each is sole writer to its SQLite)
  \item 8 concurrent readers per chain (API layer, WAL mode)
  \item Independent backup and restore per chain
  \item Drop a chain by deleting its directory
\end{itemize}

\subsection{Dual-Layer Storage}

\begin{definition}[Unified Store]
The storage layer combines two complementary backends:
\begin{itemize}
  \item \textbf{KV Layer} (BadgerDB): $O(1)$ lookups by hash, height, or composite key.
        Prefix-isolated namespaces: \texttt{blk:}, \texttt{vtx:}, \texttt{edg:}, \texttt{tx:},
        \texttt{st:}, \texttt{meta:}, \texttt{idx:}.
  \item \textbf{Query Layer} (SQLite): Complex queries with SQL, joins, aggregations,
        full-text search. WAL mode for concurrent reads.
\end{itemize}
\end{definition}

The indexer performs dual-writes to both layers. Reads prefer the KV layer for direct
lookups and fall back to the Query layer for filtered/sorted/aggregated queries.

